\scanpage{95}
\mspage{44}

est dans $W$. Je vais prouver ceci\,:

\uline{Lemme} \unsure{Soient} $S$ dans Cat, $X\in\mathrm{Ob}\,\underline{\mathrm{Fib}}_0\,S$, et considérons, pour un $S'$ dans Cat,
\[
\begin{aligned}
f&\mathrel{\unsure{\vdash|\mapsto}}f^*(X/S)=X\times_S(S',f)\\
\underline{\mathrm{Hom}}(S',S)^\circ&\longrightarrow\underline{\mathrm{Fib}}_0\,S'
\end{aligned}
\]
(Ce foncteur est défini grâce au fait que $X$ est Cat-fibrant sur $S$, et peut s'interpréter donc, \deleted{\unsure{à une}} \inserted{à $S$-}équivalence près, en termes d'un foncteur $S^\circ\mathbin{\unsure{\text{---}|\to}}\mathrm{Cat}$ \ldots) Les flèches de $\underline{\mathrm{Hom}}(S',S)$ sont transformées en flèches $\in W^d_{S'}$. \deleted{\unsure{donc}}
% slanted
\margin{\uline{NB} $\unsure{\Phi|\mathscr L}_S(X)$ était \unsure{même|justement} déjà \unsure{scindé|scindée}, \unsure{définie|défini} par le foncteur $s\mapsto{}_{s\backslash}X_{\deleted{\unsure{s}}}$}
En effet, en termes de foncteurs
\[
S'{}^\circ\longrightarrow\mathrm{Cat}
\]
si $X$ est défini par
\[
S^\circ\xrightarrow{\unsure{\mathscr X|\mathcal X}}\mathrm{Cat},
\]
alors $f^*(X)$ est défini par $\mathscr X\circ f$, et la flèche
\[
\mathscr X\circ f\mathbin{\text{---}}\mathscr X\circ g
\]
définie par $u\colon f\mathbin{\text{---}}g$ n'est autre que $\mathscr X*u$. Sa valeur en $s'\in\mathrm{Ob}\,S'$ est obtenue en prenant
\[
u(s')\colon f(s')\longrightarrow g(s')
\]
d'où par $\mathscr X$
\[
\mathscr X(u(s'))\colon
\underset{\deleted{\unsure{g}}(\mathscr X\circ g)(s')}{\tikz[baseline=(term.base)]{\node[inner sep=0pt](term){$\mathscr X(g(s'))$};\draw ([xshift=-1pt]term.south west) -- ([xshift=-1pt,yshift=-2pt]term.south west) -- ([xshift=1pt,yshift=-2pt]term.south east) -- ([xshift=1pt]term.south east);}}
\longrightarrow
\underset{(\mathscr X\circ f)(s')}{\tikz[baseline=(term.base)]{\node[inner sep=0pt](term){$\mathscr X(f(s'\unsure{)|,})$};\draw ([xshift=-1pt]term.south west) -- ([xshift=-1pt,yshift=-2pt]term.south west) -- ([xshift=1pt,yshift=-2pt]term.south east) -- ([xshift=1pt]term.south east);}}
\]
% slanted
\margin{à dégager en énoncé à part}
Mais l'hyp. \deleted{que} $X\in\underline{\mathrm{Fib}}_0(S)$ signifie que les $\mathscr X(w)$, $w\in\mathrm{Fl}(S)$, sont \unsure{ts|tous} dans $W$. \unsure{Donc|D'où}

% archive sheet number: 95
